\chapter{Inversion constitutiva and curva modular cargada}
\label{rm:chap:curva-modular-reconstruccion}

\section{The Spectral Completion from \(3\) to \(4\)}

The vacuum operator of \cref{rm:ch:vacuum} is obtained by applying to the
angular operator the function
\begin{equation}
 F(q):=\frac{1-q^{90}}{1-q^{120}},
 \qquad 0<q<1.
 \label{rm:eq:curve-Fq}
\end{equation}
The factorization of the two channels by the elementary mode \(30\),
\begin{equation}
 90=3\cdot30,
 \qquad
 120=4\cdot30,
 \label{rm:eq:curve-90-120}
\end{equation}
allows \(s=q^{30}\) to be introduced and written as
\begin{equation}
 F(q)=f(s),
 \qquad
 f(s):=\frac{1+s+s^2}{1+s+s^2+s^3}.
 \label{rm:eq:curve-fs}
\end{equation}
The numerator contains three consecutive stages of the mode, and the denominator
incorporates the fourth. The completion defect is
\begin{equation}
 1-f(s)=\frac{s^3}{1+s+s^2+s^3}>0.
 \label{rm:eq:curve-completion-defect}
\end{equation}
Thus, the fourth stage is recorded as positive memory within the
response.

\begin{theorem}[Bijectivity of the response \(3\to4\)]
\label{rm:thm:curve-F-bijection}
The function \(F\) is a decreasing analytic diffeomorphism
\begin{equation}
 \boxed{F:(0,1)\xrightarrow{\;\sim\;}(3/4,1).}
 \label{rm:eq:curve-F-range}
\end{equation}
for each \(r\in(3/4,1)\), the value \(q=F^{-1}(r)\) is obtained as
\(q=s^{1/30}\), where \(s\in(0,1)\) is the root \emph{unique} of
\begin{equation}
 r s^3+(r-1)(1+s+s^2)=0.
 \label{rm:eq:curve-inverse-cubic}
\end{equation}
\end{theorem}

\begin{proof}
The derivative with respect to \(s\) is
\begin{equation}
 f'(s)
 =-\frac{s^2(3+2s+s^2)}{(1+s+s^2+s^3)^2}<0
 \qquad(0<s<1).
 \label{rm:eq:curve-f-derivative}
\end{equation}
As \(\dd s/\dd q=30q^{29}>0\), tambiin \(F'(q)<0\). The limites are
\[
 \lim_{q\to0^+}F(q)=1,
 \qquad
 \lim_{q\to1^-}F(q)=\frac{90}{120}=\frac34.
\]
Continuity and strict monotonicity prove bijectivity; the derivative
does not vanish on the interval, so the inverse function theorem
proves analyticity. Multiplying
\(r=(1+s+s^2)/(1+s+s^2+s^3)\) by the denominator produces
\cref{rm:eq:curve-inverse-cubic}; uniqueness of its root in \((0,1)\) is
equivalent to the injectivity already proved.
\end{proof}

The inverter \cref{rm:eq:curve-inverse-cubic} provides a numerical procedure that is
intrinsic. A bisection on \((0,1)\) preserves the interval and
certifies each digit; a Newton method can accelerate it using
\cref{rm:eq:curve-f-derivative}. Decimal evaluation enters after
fixing the equation and its root, which is unique.

\section{Exact reconstruction of the angular state from the vacuum}

Let
\begin{equation}
 T=q_+P_++q_-P_-,
 \qquad
 \mathcal V=F(T)=r_+P_++r_-P_-,
 \qquad
 r_\pm=F(q_\pm).
 \label{rm:eq:curve-T-V}
\end{equation}
The two dimensionless constitutive quantities are assigned to the sheets through
\begin{equation}
 \widehat\varepsilon=r_+^2,
 \qquad
 \widehat\mu=r_-^2.
 \label{rm:eq:curve-mu-epsilon}
\end{equation}

\begin{theorem}[Reconstruction constitutive reversible]
\label{rm:thm:curve-vacuum-tomography}
The complete constitutive response determines the angular antecedent. In the
labeled spectral basis,
\begin{align}
 \boxed{q_+=F^{-1}(\sqrt{\widehat\varepsilon}),}
 &\qquad
 \boxed{q_-=F^{-1}(\sqrt{\widehat\mu}),}
 \label{rm:eq:curve-recover-qpm}\\
 \boxed{D_A=q_+q_-}
 &=F^{-1}(\sqrt{\widehat\varepsilon})
   F^{-1}(\sqrt{\widehat\mu}),
 \label{rm:eq:curve-recover-DA}\\
 \boxed{A_{\rm rad}=-\frac12\log D_A,}
 &\qquad
 \boxed{C^*_{\rm rad}=\frac12\log\frac{q_-}{q_+}.}
 \label{rm:eq:curve-recover-A-C}
\end{align}
If the operator \(\mathcal V\) is known, rather than its two labelled
eigenvalues, the reconstruction takes the form of functional calculus
\begin{equation}
 \boxed{T=F^{-1}(\mathcal V).}
 \label{rm:eq:curve-functional-inverse}
\end{equation}
\end{theorem}

\begin{proof}
The eigenvalues of \(\mathcal V\) belong to \((3/4,1)\), and the positive root of each constitutive relation is \(r_\pm\). The \cref{rm:thm:curve-F-bijection} provides a unique \(q_\pm\) for each sheet. The formulas of \(D_A,A_{\rm rad},C^*_{\rm rad}\) are the inverses of \cref{rm:eq:modular-determinant-ratio}. For the operator form, the spectral theorem applies the continuous function \(F^{-1}\) to the spectrum of \(\mathcal V\).
\end{proof}

When only the two numbers
\((\widehat\mu,\widehat\varepsilon)\) are retained, the assignment to sheets is part
of the typed data. The operator \(\mathcal V\) also contains its
spectral projectors and reconstructs the full representation. This
distinction separates spectral reconstruction from the choice of a
basis.

The remaining readings are recovered simultaneously:
\begin{equation}
 \widehat Z=\frac{r_-}{r_+}
 =\sqrt{\frac{\widehat\mu}{\widehat\varepsilon}},
 \qquad
 \widehat c=\frac1{r_-r_+}
 =\frac1{\sqrt{\widehat\mu\widehat\varepsilon}}.
 \label{rm:eq:curve-recover-Z-c}
\end{equation}
The product of the sheets publishes the common
propagation; the quotient publishes constitutive
orientation. Under \(C^*\mapsto-C^*\), the two constitutive quantities
are exchanged, \(\widehat Z\) is inverted, and \(\widehat c\) remains fixed.

\begin{corollary}[Electroweak cross-check]
\label{rm:cor:curve-weak-control}
At the angular interface where
\((m_W^{\angle}/m_Z^{\angle})^2=D_A\), the reconstruction of the vacuum
implies
\begin{equation}
 \boxed{
 \sin^2\theta_W^{\rm OS,HMT}
 =1-F^{-1}(\sqrt{\widehat\varepsilon})
      F^{-1}(\sqrt{\widehat\mu}).}
 \label{rm:eq:curve-weak-from-vacuum}
\end{equation}
\end{corollary}

\begin{proof}
The \cref{rm:thm:curve-vacuum-tomography} reconstructs \(D_A\). The stated angular
interface defines
\(\sin^2\theta_W^{\rm OS}=1-(m_W/m_Z)^2=1-D_A\).
\end{proof}

The equation \cref{rm:eq:curve-weak-from-vacuum} is a consistency check between sister
readers. The causal genealogy continues from \(T\): the vacuum
applies \(F(T)\), and the electroweak reader applies \(\det T\). The
invertibility of \(F\) permits comparison of the two publications without
making either one the foundation of the other.

\section{The loaded modular curve}

The identity
\begin{equation}
 D_A=\exp\!\left(-\frac{100\pi}{9}\alpha_{\HMT}\right)
 \label{rm:eq:curve-D-alpha}
\end{equation}
suggests using the determinant as a global coordinate of the family. For
\(D\in(0,1)\), we define
\begin{equation}
 \boxed{
 Q(D):=-\frac9{25}\log D,}
 \qquad
 a(D):=\frac{Q(D)}{4\pi}
 =-\frac9{100\pi}\log D.
 \label{rm:eq:curve-Q-a}
\end{equation}
At the HMT point, \(a(D_A)=\alpha_{\HMT}\) and
\begin{equation}
 \boxed{Q(D_A)=4\pi\alpha_{\HMT}=:e_g^2.}
 \label{rm:eq:curve-Q-gauge}
\end{equation}
The symbol \(Q(D)\) here designates the dimensionless charged coupling;
it is distinct from every dimensional section of charge or energy.

To make the regional reader self-contained, we introduce the quintic
character
\begin{equation}
 \begin{aligned}
 H_5(a):={}&\frac12\sqrt{\frac{1000a}{\varphi_{\HMT}}}-a
 +\frac9{16}a^2-\frac59a^3
 +\frac7{48}a^4-\frac1{54}a^5,
 \end{aligned}
 \label{rm:eq:curve-H5}
\end{equation}
the regional germ
\begin{equation}
 \mathcal C_\pi(a,D)
 :=169a^6+\frac{Da^7}{1-Da},
 \label{rm:eq:curve-Cpi}
\end{equation}
and the oriented phase
\begin{equation}
 y(D):=\frac{\pi}{90}\bigl(H_5(a(D))+a(D)\bigr)
 -\mathcal C_\pi(a(D),D).
 \label{rm:eq:curve-yD}
\end{equation}
The selector \(169\), the finite term of degree six, and
the geometric prolongation of degree seven retain
their regional provenance. The condition \(0<Da(D)<1\)
guarantees convergence of the germ.

We define the oriented quotient.
\begin{equation}
 \varrho_{\rm or}(D):=\ee^{2y(D)}
 \label{rm:eq:curve-rho-or}
\end{equation}
and the admissible interval
\begin{equation}
 \mathcal I_{\rm mod}:=
 \left\{D\in(0,1):
 0<Da(D)<1,
 \ |y(D)|<-\frac12\log D
 \right\}.
 \label{rm:eq:curve-domain}
\end{equation}
The second inequality ensures that the two reconstructed channels are strict contractions.

\begin{definition}[Modular curve loaded]
\label{rm:def:charged-modular-curve} The charged modular curve is the map
\begin{equation}
 \Gamma_{\rm ch}:\mathcal I_{\rm mod}\longrightarrow(0,1)\times\R_{>0},
 \qquad
 \boxed{\Gamma_{\rm ch}(D)=\bigl(D,\varrho_{\rm or}(D)\bigr).}
 \label{rm:eq:curve-Gamma}
\end{equation}
On it,
\begin{equation}
 q_-(D)=\sqrt{D\varrho_{\rm or}(D)},
 \qquad
 q_+(D)=\sqrt{\frac{D}{\varrho_{\rm or}(D)}}.
 \label{rm:eq:curve-qD}
\end{equation}
\end{definition}

\begin{theorem}[Functional rank one]
\label{rm:thm:curve-rank-one}
The curve \(\Gamma_{\rm ch}\) is an analytic immersion of rank one.
The generic space of positive bilayer contractions requires two
independent coordinates \((D,\varrho_{\rm or})\); the HMT reader restricts
that space to the graph \(\varrho_{\rm or}=\varrho_{\rm or}(D)\).
Moreover,
\begin{equation}
 A_{\rm rad}(D)=-\frac12\log D
 =\frac{50\pi}{9}a(D),
 \qquad
 C^*_{\rm rad}(D)=y(D).
 \label{rm:eq:curve-A-C-D}
\end{equation}
\end{theorem}

\begin{proof}
The functions of \cref{rm:eq:curve-Q-a,rm:eq:curve-H5,rm:eq:curve-Cpi,rm:eq:curve-yD}
are analytic on \(\mathcal I_{\rm mod}\). The first component of
\(\Gamma_{\rm ch}'(D)\) is one, so the derivative never vanishes and its rank
is one. The identities \cref{rm:eq:curve-A-C-D} follow directly from the
definitions. Finally, the product and quotient of
\cref{rm:eq:curve-qD} are \(D\) and \(\varrho_{\rm or}(D)\), respectively.
\end{proof}

The resolvent family of the chapter previous is restricted to the same curve:
\begin{equation}
 \mathfrak M_D(z)
 :=Q(D)^{-1}\bigl(zI-T(D)\bigr)^{-1},
 \qquad
 T(D)=q_+(D)P_++q_-(D)P_-.
 \label{rm:eq:curve-resolvent-D}
\end{equation}
At the point \(D_A\), the equality \(Q(D_A)=4\pi\mathfrak A_\beta\)
identifies \(\mathfrak M_D\) with \(\mathfrak M_\beta\). The coupling,
intensity, orientation, and vacuum are thus coordinated by a
single genealogical variable, while their reader functions remain
typed.

\section{Confluence between vacuum, Higgs and impedance}

The vacuum reader publishes along the curve
\begin{align}
 \widehat\mu(D)&=F(q_-(D))^2,
 &\widehat\varepsilon(D)&=F(q_+(D))^2,
 \label{rm:eq:curve-vacuum-D}\\
 \widehat Z(D)&=\frac{F(q_-(D))}{F(q_+(D))},
 &\widehat c(D)&=\frac1{F(q_-(D))F(q_+(D))}.
 \label{rm:eq:curve-Z-c-D}
\end{align}
By the injectivity of \(F\), either complete operator pair
reconstructs \((D,\varrho_{\rm or})\).

The Higgs angular route contains the orientation through
\begin{equation}
 \frac{m_H^{\angle}}{m_W^{\angle}}
 =\exp\!\left(3A_{\rm rad}+\frac53C^*_{\rm rad}\right).
 \label{rm:eq:curve-HW-ratio}
\end{equation}
In the tree interface with
\(m_H^2=2\lambda_Hv^2\) and \(m_W^2=g^2v^2/4\), this identity is equivalent to
\begin{equation}
 \lambda_H
 =\frac{Q(D)}{8(1-D)}D^{-3}
  \varrho_{\rm or}(D)^{5/3}.
 \label{rm:eq:curve-lambda-H}
\end{equation}
Therefore, the scalar
\begin{equation}
 \Xi_H(D,\lambda_H)
 :=\frac{8(1-D)D^3\lambda_H}{Q(D)}
 \label{rm:eq:curve-Xi-H}
\end{equation}
satisfies
\begin{equation}
 \boxed{
 \Xi_H=\varrho_{\rm or}^{5/3},
 \qquad
 \varrho_{\rm or}=\Xi_H^{3/5}.}
 \label{rm:eq:curve-Higgs-rho}
\end{equation}

\begin{corollary}[Higgs--impedance reconstruction]
\label{rm:cor:curve-Higgs-impedance}
Under the declared tree-level interface and for \(\Xi_H>0\),
\begin{equation}
 \boxed{
 \widehat Z
 =\frac{
 F\!\left(\sqrt D\,\Xi_H^{3/10}\right)}
 {F\!\left(\sqrt D\,\Xi_H^{-3/10}\right)}.}
 \label{rm:eq:curve-Higgs-Z}
\end{equation}
The same quotient orientado is reconstructed by three vias:
\begin{equation}
 \varrho_{\rm or}^{\rm HMT}(D)
 =\frac{F^{-1}(\sqrt{\widehat\mu})}
        {F^{-1}(\sqrt{\widehat\varepsilon})}
 =\Xi_H(D,\lambda_H)^{3/5}.
 \label{rm:eq:curve-triple-confluence}
\end{equation}
\end{corollary}

\begin{proof}
The equation \cref{rm:eq:curve-Higgs-rho} gives
\(q_-=\sqrt D\,\Xi_H^{3/10}\) and
\(q_+=\sqrt D\,\Xi_H^{-3/10}\). Substituting them in
\(\widehat Z=F(q_-)/F(q_+)\) proves \cref{rm:eq:curve-Higgs-Z}. The second
equality of \cref{rm:eq:curve-triple-confluence} is
\cref{rm:eq:curve-Higgs-rho}; the first proceeds of the
\cref{rm:thm:curve-vacuum-tomography}.
\end{proof}

The substitution of \(\lambda_H\) obtained from
\cref{rm:eq:curve-lambda-H} turns
\cref{rm:eq:curve-triple-confluence} into an algebraic consistency identity.
A determination of \(\lambda_H\) from an independent reader turns
the same equation into a falsifiable cross-check of the orientation. This
separation prevents the attribution of probative independence to a mere inversion of
the starting equation.

The gauge partition also assumes a form determined purely
by \(D\):
\begin{equation}
 g^2=\frac{Q(D)}{1-D},
 \qquad
 g'^2=\frac{Q(D)}D,
 \qquad
 \boxed{
 \frac1{g^2}+\frac1{g'^2}=\frac1{Q(D)}
 =-\frac{25}{9\log D}.}
 \label{rm:eq:curve-gauge-harmonic}
\end{equation}
The total coupling \(Q(D)\) is apportioned between the two gauge sectors,
and the harmonic sum exactly recovers the charged macro-object.

\section{Compatible causal section matching with orientation and boundary data}

The coupling \(Q(D)\) should be distinguished from the Planck energy,
which we shall denote by \(\mathcal Q_P=E_P\). The Planck pair defines an
internal velocity
\begin{equation}
 c_{\rm int}:=\frac{\mathcal Q_P\ell_P}{\hbar},
 \label{rm:eq:curve-c-int}
\end{equation}
while the constitutive operator defines
\begin{equation}
 c_{\rm vac}:=\frac1{\sqrt{\mu_{\rm vac}\varepsilon_{\rm vac}}}.
 \label{rm:eq:curve-c-vac}
\end{equation}
Both expressions live in typed causal curves. A sewing morphism
\begin{equation}
 J_C:L_C^{\rm vac}\longrightarrow L_C^{\rm grav}
 \label{rm:eq:curve-JC}
\end{equation}
allows their comparison, and their dimensionless defect is
\begin{equation}
 \boxed{
 \Xi_c
 :=\frac{\mathcal Q_P\ell_P}{\hbar}
   \sqrt{\mu_{\rm vac}\varepsilon_{\rm vac}}
 =\frac{c_{\rm int}}{c_{\rm vac}}.}
 \label{rm:eq:curve-Xi-c}
\end{equation}

\begin{criterion}[Causal seam]
\label{rm:crit:curve-causal-seam}
The electromagnetic vacuum and the gravitational realization publish the
same causal section exactly when
\begin{equation}
 \boxed{\Xi_c=1,}
 \qquad
 J_C(c_{\rm vac})=c_{\rm int}.
 \label{rm:eq:curve-Xi-one}
\end{equation}
\end{criterion}

Once this seam is satisfied, the CT108 identity
\(
G=(54/\pi)^2c^5t_0^2/\hbar
\)
is expressed entirely through the constitutive product:
\begin{equation}
 \boxed{
 G=\left(\frac{54}{\pi}\right)^2
 \frac{t_0^2}
 {\hbar(\mu_{\rm vac}\varepsilon_{\rm vac})^{5/2}}.}
 \label{rm:eq:curve-G-from-vacuum}
\end{equation}
Gravity uses the even product of the leaves; impedance and charge
use the oriented quotient
\begin{equation}
 Z_{\rm vac}^2=\frac{\mu_{\rm vac}}{\varepsilon_{\rm vac}}.
 \label{rm:eq:curve-Z-oriented}
\end{equation}
The pair product--quotient distributes, therefore, causality and orientation
without disaggregating the spectral antecedent.

\section{Stable reconstruction in the UHF tower}

Let \(\tau_m=\tau_2\otimes\tau_{9^m}\) be the normalized trace of
\(\mathcal A_m=M_2\otimes M_{9^m}\), and let
\begin{equation}
 T_m=T\otimes I_{9^m},
 \qquad
 \mathcal V_m=F(T_m)=\mathcal V\otimes I_{9^m}.
 \label{rm:eq:curve-UHF-TV}
\end{equation}
Functional inversion commutes with refinement:
\begin{equation}
 F^{-1}(\mathcal V_m)=T_m,
 \qquad
 E_m(\mathcal V_{m+1})=\mathcal V_m,
 \qquad
 E_m(T_{m+1})=T_m.
 \label{rm:eq:curve-UHF-inversion}
\end{equation}

The ordinary matrix determinant of \(T_m\) is \(D_A^{9^m}\). The
invariant compatible with refinement is the normalized logarithmic
determinant
\begin{equation}
 \boxed{
 \mathfrak d_m(T_m)
 :=\exp\!\left(2\tau_m(\log T_m)\right)=D_A.}
 \label{rm:eq:curve-normalized-determinant}
\end{equation}
The orientation admits the equally stable prolongation
\begin{equation}
 \boxed{
 \mathfrak r_m(T_m,R_m)
 :=\exp\!\left(-2\tau_m(R_m\log T_m)\right)
 =\varrho_{\rm or},}
 \qquad
 R_m=R\otimes I_{9^m}.
 \label{rm:eq:curve-normalized-orientation}
\end{equation}

\begin{proposition}[Stability of reconstruction at every depth]
\label{rm:prop:curve-UHF-tomography}
For every \(m\), the pair
\begin{equation}
 \left(
 \mathfrak d_m(F^{-1}(\mathcal V_m)),
 \mathfrak r_m(F^{-1}(\mathcal V_m),R_m)
 \right)
 \label{rm:eq:curve-UHF-pair}
\end{equation}
equals \((D_A,\varrho_{\rm or})\). Consequently, it reconstructs the same coordinates \(A_{\rm rad}\) and \(C^*_{\rm rad}\) at every level and in the UHF limit.
\end{proposition}

\begin{proof}
Since
\(\log T_m=(\log T)\otimes I_{9^m}\), the normalized trace eliminates the
amplification factor. Moreover,
\[
 \tau_2(\log T)=\frac12(\log q_++\log q_-)=\frac12\log D_A
\]
and
\[
 \tau_2(R\log T)=-y=-\frac12\log\varrho_{\rm or}.
\]
The equations \cref{rm:eq:curve-normalized-determinant,rm:eq:curve-normalized-orientation}
follow upon exponentiation. Functional inversion of
\cref{rm:eq:curve-UHF-inversion} completes the reconstruction.
\end{proof}

This formulation shows that the completion \(3\to4\) preserves its
reconstructive capacity at every depth. The nonadic cylinders
increase the resolution of the state; the normalized determinant,
orientation, response, and modular curve remain natural with respect
to the expectations.

\section{Separation between internal quantum identities, physical interfaces, and reconstruction obstructions}

\begin{proofstatusbox}
The monotonicity and bijectivity of \(F\), the operatorial inversion, the
reconstruction \((q_+,q_-,D_A,C^*)\), the rank one of
\(\Gamma_{\rm ch}\), the harmonic gauge identity, and UHF stability
are exact results. The charged modular curve formalizes a pre-existing
authorial architecture with explicit starting structure: closure of
\(\alpha_{\HMT}\), regional selector \(169\), character \(H_5\), germ
\(\mathcal C_\pi\), and orientation. The
Higgs--impedance equality is exact within the tree interface; it constitutes an
independent proof only when \(\lambda_H\) follows from a distinct reader.
The condition \(\Xi_c=1\) and the subsequent gravitational equation
are a typed physical interface whose promotion requires the morphism \(J_C\).
\end{proofstatusbox}

\begin{falsifierbox}
The reconstruction is refuted if \(F'(q)\) changes sign in \((0,1)\), if the cubic equation admits two roots in that interval, if \(F^{-1}(\mathcal V)\) does not recover \(T\), or if the constitutive relations reconstruct a product different from \(D_A\). The curve is refuted if \(Q(D_A)\ne4\pi\alpha_{\HMT}\), if its derivative loses rank, if \(q_\pm(D)\) leaves \((0,1)\) within the declared domain, or if the three reconstructions of \(\varrho_{\rm or}\) in \cref{rm:eq:curve-triple-confluence} disagree under the same normalization. The causal seam is refuted by \(\Xi_c\ne1\) after applying the line morphism. In the tower, using the ordinary determinant as though it were invariant would produce \(D_A^{9^m}\); the correct control requires \cref{rm:eq:curve-normalized-determinant}.
\end{falsifierbox}
